\section{Recovery of the original hypersurface and the converse}
\label{q4:rec:section}

The static classification identifies a metric on the orbit space.  We now
recover the original spacelike hypersurface, including its second fundamental
form and its smooth structure at the horizon.  We then compute the invariant
ADM mass of every admissible Tangherlini slice, allowing a nonzero asymptotic
boost.

\subsection{The original graph in the open exterior}

Throughout the forward implication, assume equality and all the horizon
hypotheses of Definition~\ref{q4:def:horizon}.  Retain
\[
 m=\sqrt{E^2-|P|^2},\qquad r_0=\sqrt{2m},\qquad
 \kappa=r_0^{-1}.
\]
Proposition~\ref{q4:cla:static-classification} identifies the whole open base
$\operatorname{int}\Omega$ with the Tangherlini spatial exterior of mass $m$:
\begin{equation}
 h_b=N^{-1}\dd r^2+r^2g_{\mathbb S^3},\qquad
 N=1-\frac{r_0^2}{r^2},\qquad r>r_0.
 \label{q4:rec:classified-base}
\end{equation}
In particular $N>0$ throughout the original interior, the open base is simply
connected, and its regular boundary attachment is homeomorphic to the original
attachment of $S$.  The adjoint field is the original field of
Theorem~\ref{q4:adj:causal-field}; all occurrences of $X^\flat$ below use the
original metric $g$.

\begin{lemma}[Recovery on the open exterior]
\label{q4:rec:interior-graph}
There is a smooth function $T_0$ on $\operatorname{int}\Omega$ such that the
map
\[
 x\longmapsto\bigl(T_0(x),r(x),\vartheta(x)\bigr)
\]
into the static Tangherlini exterior induces the original $g$ and $K$.
Here $(r,\vartheta)$ are the global base coordinates in
Equation~\eqref{q4:rec:classified-base}, with $\vartheta\in\mathbb S^3$.
\end{lemma}

\begin{proof}
The one-form $A_b=N^{-1}X^\flat$ is closed by the staticity conclusion.  Simple
connectedness gives a global primitive with
\begin{equation}
 \dd T_0=-A_b.
 \label{q4:rec:primitive}
\end{equation}
On $\mathbb R\times\operatorname{int}\Omega$ put $T=t+T_0(x)$.  The definitions
of $h_b$, $A_b$, and $N=u^2-|X|_g^2$ give the exact identity
\begin{equation}
 \begin{split}
 -N\dd T^2+h_b
 &=-N(\dd t-A_b)^2+h_b\\
 &=-u^2\dd t^2+
   g_{ij}(\dd x^i+X^i\dd t)(\dd x^j+X^j\dd t).
 \end{split}
 \label{q4:rec:metric-identity}
\end{equation}
Thus the graph $T=T_0$ is exactly the original $t=0$ slice and induces $g$.
Static time is future increasing, and its future unit normal is
$n=u^{-1}(\partial_t-X)$.  With the convention of the theorem, the
lapse--shift identity is
\[
 \partial_tg=2uK+\mathcal L_Xg.
\]
The metric on the right of Equation~\eqref{q4:rec:metric-identity} is stationary,
and Equation~\eqref{q4:adj:kid-first} therefore gives its second form as
$-\mathcal L_Xg/(2u)=K$.  This proves the claim for the original tensor,
without a further deformation.  These local sign conventions also agree
with the Killing initial data identities
in~\cite[Equations~(2.6) and~(2.14)]{BeigChrusciel1997KID}.
\end{proof}

\subsection{Extension through the horizon in the original smooth structure}

The base has two possible regular boundary structures.  When $u|_S>0$,
$N$ is an original smooth defining function, whereas the regular base
coordinate is $\lambda=\sqrt N$.  When $u|_S=0$, the base is smooth in the
original collar coordinate.  The following lemma treats these cases before
passing to regular spacetime coordinates.

\begin{lemma}[Smooth horizon extension]
\label{q4:rec:horizon-extension}
The graph in Lemma~\ref{q4:rec:interior-graph} extends smoothly, in the original
smooth structure of $\Omega$, to $S$ in the regular maximal Tangherlini
extension.  Its boundary is a smooth section of the corresponding future
horizon if $u|_S>0$, and is the bifurcation sphere if $u|_S=0$.
\end{lemma}

\begin{proof}
Connectedness of $S$ and Theorem~\ref{q4:adj:causal-field} leave exactly the two
cases just described.

\paragraph{Positive boundary lapse.}
Equation~\eqref{q4:adj:positive-collar} states that
$\dd N=2\kappa X^\flat$ on $S$ and that $N$ is an original defining function.
Consequently the smooth covector
$X^\flat-(2\kappa)^{-1}\dd N$ vanishes on $S$.  Dividing its components by $N$
in an original collar gives a smooth one-form $\beta_0$ such that
\begin{equation}
 A_b=\frac1{2\kappa}\dd\log N+\beta_0,\qquad
 T_0=-\frac1{2\kappa}\log N+B_0.
 \label{q4:rec:log-time}
\end{equation}
Here $B_0$ is smooth up to $S$: its differential is $-\beta_0$, and its values
on an interior collar section determine its smooth extension by integration
along collar segments.  The primitive in
Equation~\eqref{q4:rec:primitive} is already single-valued, so this construction
has no period ambiguity.

We also need the angular coordinate $\vartheta$ to be smooth in the original
coordinate $N$, not only in $\lambda$.  The specific reflection matters here.
In the original-derived coordinates $(\lambda,y)$, with $N=\lambda^2$,
Equation~\eqref{q4:sta:positive-lapse-base-collar} has even $\lambda\lambda$ and
$AB$ coefficients and odd $\lambda A$ coefficients, all original functions
being evaluated at $(\lambda^2,y)$.  Thus the smooth isometric reflection is
exactly
\[
 \mathcal R(\lambda,y)=(-\lambda,y).
\]
Let $p$ be normal projection to $S$ in this reflected base collar.
Reflection fixes $S$ pointwise and reverses its unit normals, so
$p\circ\mathcal R=p$.  The smooth one-sided isometry in
Proposition~\ref{q4:cla:static-classification} sends these normal geodesics to
the radial normal geodesics of the classified base.  If $\vartheta_S$ denotes
its smooth angular boundary identification, the original angular map on the
positive collar is therefore $\vartheta=\vartheta_S\circ p$.  This formula
extends it smoothly to the reflected collar and makes it invariant under
the displayed $\mathcal R$.  In particular it is even in the specific
$\lambda$ coordinate obtained from original $N$, with its tangential
parameters fixed.  A smooth even function of
$\lambda$ is a smooth one-sided function of $\lambda^2$: Taylor expansion
has only even powers, and the differentiated remainder gives this assertion
to every finite order.  Applying this in angular charts proves that
$\vartheta$ is smooth in the original $(N,y)$ collar.

\paragraph{Zero boundary lapse.}
Let $s\geq0$ be original outward collar distance.  By
Equation~\eqref{q4:adj:zero-collar}, there are smooth $d,Y_2$ with
\begin{equation}
 u=sd,\qquad X=s^2Y_2,\qquad d|_S=\kappa,\qquad
 N=s^2\bigl(d^2-s^2|Y_2|_g^2\bigr).
 \label{q4:rec:zero-collar}
\end{equation}
It follows that $A_b$, and hence $T_0$, extend smoothly and finitely to $S$.
The positive square root on the exterior is
\[
 \lambda=s\sqrt{d^2-s^2|Y_2|_g^2},
\]
which is smooth one-sided in $s$.  In this case the original and regular
one-sided base structures agree.  The smooth one-sided base identification
therefore makes $\vartheta$ smooth in the original collar as well.

\paragraph{Regular spacetime coordinates.}
Define the tortoise coordinate and Kruskal coordinates by
\begin{equation}
 \begin{split}
 r_*&=r+\frac{r_0}{2}\log\frac{r-r_0}{r+r_0},\\
 \mathsf U&=-\exp\bigl(-\kappa(T-r_*)\bigr),\qquad
 \mathsf V=\exp\bigl(\kappa(T+r_*)\bigr).
 \end{split}
 \label{q4:rec:kruskal}
\end{equation}
Differentiation gives $\dd r_*/\dd r=N^{-1}$.  More precisely,
\begin{equation}
 e^{\kappa r_*}=\sqrt N\,D_*(N),\qquad
 D_*(N)=e^{r/r_0}\frac{r}{r+r_0},\qquad
 r=\frac{r_0}{\sqrt{1-N}}.
 \label{q4:rec:regular-factor}
\end{equation}
Thus $D_*$ is smooth positive near $N=0$, with $D_*(0)=e/2$.
The two-dimensional static part becomes
\[
 -N\dd T^2+N^{-1}\dd r^2
   =-\frac{1}{\kappa^2D_*(N)^2}\,\dd\mathsf U\,\dd\mathsf V.
\]
Moreover $-\mathsf U\mathsf V=N D_*(N)^2$ has a smooth local inverse for
$N$ near zero.  These are regular horizon coordinates: the coefficient of
$\dd\mathsf U\,\dd\mathsf V$ is smooth negative nonzero and $r$ is smooth
across the horizon.

In the positive boundary-lapse case, substituting
Equation~\eqref{q4:rec:log-time} into Equation~\eqref{q4:rec:kruskal} on the graph
gives
\begin{equation}
 \mathsf U=-N D_*(N)e^{-\kappa B_0},\qquad
 \mathsf V=D_*(N)e^{\kappa B_0}.
 \label{q4:rec:future-section}
\end{equation}
Both functions are smooth in the original collar, and $S$ maps to
$\mathsf U=0$, $\mathsf V>0$.  Since its angular map identifies $S$ with
$\mathbb S^3$, this is a smooth section of the future horizon.

In the zero boundary-lapse case, the same graph formulas are
\[
 \mathsf U=-\lambda D_*(N)e^{-\kappa T_0},\qquad
 \mathsf V=\lambda D_*(N)e^{\kappa T_0}.
\]
Equation~\eqref{q4:rec:zero-collar} makes them smooth in $s$, and both vanish on
$S$.  Together with the angular identification, they map $S$ onto the
bifurcation sphere.  This proves the asserted smooth extension in both cases.
\end{proof}

\begin{proposition}[Global recovery of the original exterior]
\label{q4:rec:embedding}
Equality under the horizon hypotheses gives a global smooth spacelike
embedding of the original $\Omega$, including $S$, in the regular maximal
Tangherlini extension of mass $m$.  The embedding pulls back the spacetime
metric and the future second fundamental form to the original $g$ and $K$.
Its image lies in the closure of the corresponding exterior, its boundary
has the form stated in Lemma~\ref{q4:rec:horizon-extension}, and its end
approaches the corresponding spatial infinity.
\end{proposition}

\begin{proof}
Combine the maps in Lemmas~\ref{q4:rec:interior-graph}
and~\ref{q4:rec:horizon-extension}.  Their pullback metric is $g$ on the open
exterior and remains $g$ at $S$ by continuity.  Since $g$ is positive definite,
the extended map is a spacelike immersion there.

It is injective on the interior by the global base identification, and on the
boundary by the round-sphere identification.  No interior point can have the
same image as a boundary point, since their area radii are respectively
$r>r_0$ and $r=r_0$.  It is also proper.  A compact set of the regular
spacetime has bounded area radius; its preimage is a closed subset of a
bounded annulus in the attached base.  Such an annulus is compact, and the
base attachment is homeomorphic to the original attachment.  A proper
injective immersion is an embedding.

To obtain the normal at the boundary, take a smooth future timelike field in
regular spacetime coordinates, project it to the normal line of the spacelike
image, and normalize.  Its normal projection is timelike and nonzero; the
result is a smooth future unit normal.  It agrees with the interior normal
by uniqueness of the future choice.  Hence the second-form identity from
Lemma~\ref{q4:rec:interior-graph} extends to $S$.  If desired, the smooth collar
expressions extend locally across $S$ as a spacelike immersion, since
immersion and spacelikeness are open conditions.  No data beyond $S$ are
prescribed or asserted.

Finally the inverse-base identity gives
\begin{equation}
 |\dd T_0|_{h_b}^2
 =\frac{|X|_g^2}{u^2N}
 \longrightarrow \frac{|b|^2}{(b^0)^2}<1,
 \label{q4:rec:spatial-infinity}
\end{equation}
using Equation~\eqref{q4:adj:decay} and $N\to1$.
Choose $a<1$ slightly larger than the limiting slope.  Increasing a fixed
base radius if necessary absorbs the radial length factor $N^{-1/2}$, so
integration along base radial rays, starting on a compact sphere, gives
$|T_0|\leq ar+C$ uniformly on the end.  Thus the end remains in a spacelike
cone of the static exterior and approaches its spatial infinity.
\end{proof}

\subsection{Asymptotic geometry of an admissible Tangherlini slice}

For the converse write $M_0>0$ for the geometric mass parameter of the
ambient Tangherlini spacetime, to distinguish it from the invariant ADM mass
that we will compute.  In isotropic spatial coordinates $y$ on its chosen
exterior, with future static time $T$, its metric is
\begin{equation}
 \mathbf g_{M_0}
 =-\left(\frac{1-M_0/(2|y|^2)}{1+M_0/(2|y|^2)}\right)^2\dd T^2
   +\left(1+\frac{M_0}{2|y|^2}\right)^2
       \sum_{i=1}^4(\dd y^i)^2.
 \label{q4:rec:isotropic-spacetime}
\end{equation}
The following argument applies to the end of every slice in the converse
class.  It derives its asymptotic plane from the prescribed intrinsic decay.

\begin{lemma}[Asymptotic affine plane]
\label{q4:rec:asymptotic-plane}
Let a smooth spacelike hypersurface in the Tangherlini exterior of mass
$M_0>0$ have an end with coordinates $x\in\mathbb R^4\setminus\overline B$
and induced data satisfying
\[
 g-\delta=O_2(|x|^{-q}),\qquad K=O_1(|x|^{-1-q}),\qquad q>1.
\]
Assume the hypersurface is smoothly embedded with a compact remaining part
and boundary, if present, in the exterior closure.  For any
$1<q_0<\min(q,2)$ there are constants
$L\in\operatorname{GL}(4,\mathbb R)$, $a\in\mathbb R^4$,
$L_0\in\mathbb R^4$, and $a_0\in\mathbb R$ such that on the end
\begin{equation}
 \begin{split}
 y&=Lx+L_0+O_2(|x|^{1-q_0}),\\
 T&=a\cdot x+a_0+O_2(|x|^{1-q_0}),\qquad
 L^tL-a\otimes a=I.
 \end{split}
 \label{q4:rec:affine-plane}
\end{equation}
In particular the limiting tangent plane is spacelike and $x$ gives
Euclidean orthonormal coordinates on it.
\end{lemma}

\begin{proof}
We first control the spatial projection and then obtain its affine asymptotics.

\paragraph{The end leaves every bounded area radius.}
Gauss and Codazzi express the all-tangential and one-normal spacetime
curvature components in a slice orthonormal frame in terms of
$\operatorname{Rm}_g$, $K^2$, and $\nabla K$.  All tend to zero by the stated
decay.  Vacuum $\mathbf{Ric}_{ij}=0$ expresses the two-normal curvature
components as sums of all-tangential ones.  Thus every component tends to
zero, and so does the spacetime curvature-square invariant.

In the static orthonormal frame of Tangherlini, curvature components with
one time index and three spatial indices vanish.  All remaining squared
components contribute nonnegatively to that invariant.  The angular
sectional curvature, computed from the static warped product, is
\[
 \frac{1-|\nabla r|_{h_{\mathrm{static}}}^2}{r^2}
 =\frac{2M_0}{r^4}.
\]
The invariant is consequently bounded below by a positive constant times
$M_0^2/r^8$, including at the horizon by regularity.  Hence the area radius
$r$, and therefore $|y|$, tend to infinity as $|x|\to\infty$.

\paragraph{A finite timelike limiting normal.}
For the estimates in original end coordinates, write $R=|x|$.
Decompose the future static Killing field along the slice as
$\partial_T=u n+X$.  On its exterior tail $u>|X|_g$.  Local translates along
this transverse field give stationary lapse--shift coordinates.  The vacuum
Killing initial data equations, with the present second-form convention, are
\begin{equation}
 \nabla_{(i}X_{j)}=-uK_{ij},\qquad
 \nabla^2u=u(\Ric_g+\tau K-2K^2)-\mathcal L_XK.
 \label{q4:rec:vacuum-kid}
\end{equation}
These follow respectively from the Killing equation and its normal
derivative using the vacuum spatial Ricci equation.  Their local form is
dimension independent; we use them here in four spatial dimensions.  The
same local identities are recorded in
\cite[Section~2]{BeigChrusciel1997KID}.

Commuting derivatives in the first equation expresses $\nabla^2X$ as
curvature times $X$ and permutations of $\nabla(uK)$.  If $Y$ denotes the
coordinate components of the pair $(u,X)$, the prescribed decay therefore
gives
\begin{equation}
 |\partial^2Y|
 \leq C|x|^{-1-q_0}|\partial Y|+C|x|^{-2-q_0}|Y|.
 \label{q4:rec:jet-estimate}
\end{equation}
Here only two derivatives of $g$ and one of $K$ are used.

The metric $g$ is uniformly comparable to the Euclidean metric on the end,
the lapse and shift are causal, and $q_0>1$.  Applying
Lemma~\ref{q4:adj:causal-jet-asymptotics} to
Equation~\eqref{q4:rec:jet-estimate} gives finite constant limits with
$O_2(R^{-q_0})$ errors.  Their normalization follows separately from
$u^2-|X|_g^2=-\mathbf g_{M_0}(\partial_T,\partial_T)\to1$,
using the already proved $|y|\to\infty$.  Thus
\begin{equation}
 u=u_\infty+O_2(R^{-q_0}),\qquad
 X=X_\infty+O_2(R^{-q_0}),\qquad
 u_\infty^2-|X_\infty|_\delta^2=1.
 \label{q4:rec:limiting-normal}
\end{equation}

\paragraph{Proper spatial projection and affine jets.}
Set $N=u^2-|X|_g^2$.  Pairing $\partial_T$ with tangent vectors to the slice
gives the exact identities
\begin{equation}
 \dd T=-\frac{X^\flat}{N},\qquad
 y^*h_{\mathrm{static}}
   =h_b=g+N^{-1}X^\flat\otimes X^\flat.
 \label{q4:rec:projection-identities}
\end{equation}
The spatial projection $y$ is a local diffeomorphism, since the static
Killing field is timelike and transverse to the spacelike hypersurface.
Equation~\eqref{q4:rec:limiting-normal} makes $h_b$ asymptotic to a positive
constant matrix.  Because $|y|\to\infty$, both its Jacobian and inverse
Jacobian are uniformly bounded far out.

We make the target coverage and path lifting explicit.  Choose $R_*$ large
enough that the preceding Jacobian bounds hold on the original end
$E_*:=\{|x|>R_*\}$.  Choose $R_0$ larger than the maximum of $|y|$ on
$\partial E_*$, and put
\[
 V_0=\{y\in\mathbb R^4:|y|>R_0\},\qquad
 U_0=E_*\cap y^{-1}(V_0).
\]
The local diffeomorphism $y:U_0\to V_0$ is proper.  Indeed a sequence whose
images remain in a compact subset of $V_0$ cannot escape to original
infinity, by $|y|\to\infty$, or approach $\partial E_*$, by the choice of
$R_0$.  It therefore has a convergent subsequence in $U_0$.  The image is
open by local invertibility and closed by properness.  It is nonempty, and
$V_0$ is connected, so the image is all of $V_0$.

A proper local diffeomorphism is a covering: inverse neighborhoods of its
finite fibers give evenly covered neighborhoods, with properness excluding
additional sheets arriving from outside those neighborhoods.  In particular,
radial paths down to a fixed sphere $|y|=R_1>R_0$ lift.  Their endpoints lie
in the compact preimage of that sphere.  The uniform inverse-Jacobian bound
gives a lifted length at most $C(|y|-R_1)$ and hence
$|x|\leq C'(1+|y|)$.  Conversely, integrating the Jacobian bound along
original coordinate rays gives
$|y|\leq C(1+|x|)$.  Hence $|y|\asymp|x|$.

The connection transformation law for
Equation~\eqref{q4:rec:projection-identities} now yields
$\partial^2y=O(|x|^{-1-q_0})$: the connection of $h_b$ has that order, while
the static connection is $O(|y|^{-3})$, and $q_0<2$.  The first identity in
Equation~\eqref{q4:rec:projection-identities} similarly gives the required
Hessian bound for $T$.  Integration on radial rays and comparison on spheres
first yield constant gradient limits and then constant translation limits.
The latter comparison costs $O(R^{1-q_0})$, which tends to zero.  This proves
the two expansions in Equation~\eqref{q4:rec:affine-plane}.  Finally the limit
of the induced metric is $L^tL-a\otimes a=I$.  In particular $L$ is
invertible and the affine plane is spacelike.
\end{proof}

The preceding lemma has recovered a genuine spacelike affine plane at
infinity from the intrinsic weak decay, rather than assuming a preferred
rest slicing. We can therefore compute the charges on that plane and show
that the decaying displacement of the actual slice changes no ADM flux.

\subsection{The invariant ADM mass of the slice}

\begin{proposition}[Mass of every admissible Tangherlini slice]
\label{q4:rec:converse-mass}
Every hypersurface in the converse class of Theorem~\ref{q4:thm:rigidity}, in
Tangherlini spacetime of geometric mass $M_0>0$, has invariant ADM mass $M_0$.
More precisely, in coordinates adapted to its asymptotic boost, its charges
are $E=M_0\Gamma$ and $P=M_0\Gamma v\,e_1$, where
$0\leq v<1$ and $\Gamma=(1-v^2)^{-1/2}$.
\end{proposition}

\begin{proof}
We compare the slice with its limiting plane and compute the flux on that plane.

\paragraph{Comparison with the affine plane.}
Apply Lemma~\ref{q4:rec:asymptotic-plane}.  Use the same parameter $x$ on the
affine plane in Equation~\eqref{q4:rec:affine-plane}, and denote its induced data
by $(g^*,K^*)$.  In affine Minkowski coordinates adapted to this plane, the
residual displacement of the actual embedding is a spacetime vector $Z$,
with time component $Z^0$, satisfying $Z=O_2(|x|^{1-q_0})$.  The ambient
metric differs from Minkowski by $O_2(|x|^{-2})$ near both embeddings.
Consequently
\begin{equation}
 \begin{split}
 g_{ij}-g^*_{ij}
   &=\partial_iZ_j+\partial_jZ_i+O_1(|x|^{-2q_0}),\\
 K_{ij}-K^*_{ij}
   &=\partial_i\partial_jZ^0+O(|x|^{-1-2q_0}).
 \end{split}
 \label{q4:rec:gauge-comparison}
\end{equation}
Here the indices on the spatial components of $Z$ are Euclidean.
To check the error orders, the quadratic Minkowski metric term is
$O_1(|x|^{-2q_0})$.  Evaluation of the ambient perturbation at the displaced
point, or its contraction with a displaced tangent vector, gives
$O_1(|x|^{-2-q_0})$, which is smaller since $q_0<2$.  For the second form use
$-\langle n,\nabla_i\partial_j\iota\rangle$.  The normal differs from the
constant future affine normal by $O(|x|^{-q_0})$; multiplying this by
$\partial^2Z$ gives $O(|x|^{-1-2q_0})$.  Ambient connection and evaluation
errors have order $O(|x|^{-3-q_0})$, again smaller.  Euclidean trace reversal
may be used in the difference with errors of the same allowed order.

The only possible additional limiting fluxes in
Equation~\eqref{q4:rec:gauge-comparison} come from its displayed linear terms.
For the metric term the energy vector is
\[
 F_i=\Delta Z_i-\partial_i\operatorname{div}Z,
 \qquad \partial_iF_i=0.
\]
For the second-form term the momentum tensor is
\[
 B_{ij}=\partial_i\partial_jZ^0-\delta_{ij}\Delta Z^0,
 \qquad \partial_jB_{ij}=0.
\]
Extend $Z$ smoothly by a cutoff through the bounded Euclidean interior.
The divergence theorem makes both linear fluxes exactly zero on every
sufficiently large sphere.  The remaining flux errors are
$O(R^{2-2q_0})$, and hence tend to zero because $q_0>1$.  It is therefore
enough to compute the charges of the affine plane in the ambient metric
\eqref{q4:rec:isotropic-spacetime}.

\paragraph{The model-plane charges.}
Translations change the ambient leading mass terms only by $O(|x|^{-3})$
and do not affect the charges. After rotations and an orthonormal
choice of $x$, the plane is the one in Lemma~\ref{q:converse:lem:boost}
with $n=4$, $p=2$, $k=3$, $M=M_0$, and $\gamma=\Gamma$.
Its static and boosted times are the present $T$ and $t$, respectively;
the ambient metric and future-normal convention agree.

Put $s_*^2=\Gamma^2x_1^2+|x_\perp|^2$. Specializing
Equations~\eqref{q:converse:eq:boost-energy-vector} and
\eqref{q:converse:eq:boost-momentum-tensor} gives
\begin{equation}
 \partial_jg^*_{ij}-\partial_ig^*_{jj}
   =6M_0\Gamma^2x_i s_*^{-4}+o(|x|^{-3}),
 \label{q4:rec:boost-energy-integrand}
\end{equation}
and
\begin{equation}
 K^*-(\tr_{g^*}K^*)g^*
 =3M_0\Gamma^2v s_*^{-4}
 \begin{pmatrix}
  x_1&x_\perp^t\\
  x_\perp&-\Gamma^2x_1I_3
 \end{pmatrix}
 +o(|x|^{-3}).
 \label{q4:rec:boost-momentum-integrand}
\end{equation}
The model metric and tensor remainders are respectively
$O_d(|x|^{-4})$ and $O_d(|x|^{-5})$, with vanishing ADM fluxes.
Equation~\eqref{q:converse:eq:ellipsoid-integral} becomes
\begin{equation}
 \int_{\mathbb S^3}
   \bigl(\Gamma^2n_1^2+|n_\perp|^2\bigr)^{-2}\dd A
 =\frac{\omega}{\Gamma}.
 \label{q4:rec:ellipsoid-integral}
\end{equation}
The energy and momentum factors are
$1/(2k\omega)=1/(6\omega)$ and $1/(k\omega)=1/(3\omega)$.
The same lemma therefore gives
\begin{equation}
 E=M_0\Gamma,\qquad P_1=M_0\Gamma v,
 \qquad P_A=0\quad(A=2,3,4).
 \label{q4:rec:boost-charges}
\end{equation}
The preceding weak-decay comparison transfers these charges to the
original slice, proving $\sqrt{E^2-|P|^2}=M_0$.
\end{proof}

\begin{lemma}[Volume of a horizon section]
\label{q4:rec:horizon-area}
A smooth spacelike section of a future Tangherlini horizon of geometric
mass $M_0>0$, including its bifurcation sphere, has induced metric
$(2M_0)g_{\mathbb S^3}$ and three-volume
$\omega(2M_0)^{3/2}$.
\end{lemma}

\begin{proof}
In the regular coordinates of Equation~\eqref{q4:rec:kruskal}, the future
horizon is $\mathsf U=0$ with area radius $r=\sqrt{2M_0}$.  The generator
direction $\partial_{\mathsf V}$ is the kernel of its induced form.  Thus
pullback to any section leaves precisely $r^2g_{\mathbb S^3}$, regardless
of the value of $\mathsf V$ along the section.  A compact connected spacelike
section has an angular projection which is a local diffeomorphism and hence
a covering of $\mathbb S^3$; compactness gives surjectivity, and simple
connectedness of $\mathbb S^3$ makes that covering a diffeomorphism.  The
same induced metric holds at the bifurcation sphere.  Taking its volume
proves the formula.
\end{proof}

\begin{proof}[Completion of the proof of Theorem~\ref{q4:thm:rigidity}]
For equality data, Theorem~\ref{q4:thm:numerical-provider} gives $m>0$.
The preceding adjoint, staticity, and classification results apply to the
original exterior, and Proposition~\ref{q4:rec:embedding} supplies the required
global smooth embedding with the original $g$ and $K$, including its normal
and its horizon boundary.  The ambient mass is the same
$m=\sqrt{E^2-|P|^2}$ by the scale in
Equation~\eqref{q4:rec:classified-base}.

Conversely, take a Tangherlini hypersurface satisfying every exterior,
decay, and horizon hypothesis in the theorem.  If its ambient geometric
mass is $M_0$, Proposition~\ref{q4:rec:converse-mass} gives invariant ADM mass
$M_0$, and Lemma~\ref{q4:rec:horizon-area} gives
$A=\omega(2M_0)^{3/2}$.  Outer area minimization, which is part of this
converse class, identifies $A$ with the original enclosing infimum.
Therefore
\[
 \sqrt{E^2-|P|^2}=M_0
   =\frac12\left(\frac A\omega\right)^{2/3}.
\]
This proves both directions, including slices with nonzero original $K$
or nonzero original ADM momentum.
\end{proof}
